\section{Polarisation dodécaphasique, dualité réticulaire et réciprocité thêta}
\label{sec:leech-dualidad}

L'incidence exceptionnelle et la polarisation du registre dodécaphasique proviennent de deux réalisations du même état enrichi. La première conserve les relations de frontière, le code de recollement et l'origine marquée qui déterminent le réseau ; la seconde fait agir une représentation sur les douze coordonnées entières déjà obtenues par les lecteurs du Topological Phase Kernel. Leur composition permet d'étudier une déformation métrique dont le paramètre modifie les longueurs relatives de douze plans et conserve simultanément le volume total et une symétrie ternaire explicite.

Le point de départ causal demeure la construction APP--TRIT--TPK de l'état enrichi et du continu conjoint. Le registre \(K\) est reçu après le registre dodécaphasique et ses opérations d'incidence et d'orientation. La représentation employée agit sur cette sortie et conserve ses douze valeurs comme des données déjà déterminées. La réalisation réticulaire est reçue avec sa forme bilinéaire, ses classes de recollement et son vecteur marqué. Ces structures, développées respectivement dans \cite{hmtX,HMT-IV}, sont maintenues lors de leur composition. La collection construite ci-dessous est une réalisation géométrique de ces données avec une coordinatisation commune à douze positions.

\subsection{Polarisation du registre et marquage des douze plans}
\label{subsec:leech-polarizacion}

Le registre dodécaphasique et sa somme sont
\begin{equation}
 \begin{split}
 K={}&(234,543,140,729,659,824,\\
     &\hspace{13mm}621,58,914,794,146,601),\\
 \sum_{j=1}^{12}K_j={}&6263.
 \end{split}
 \label{eq:leech-K}
\end{equation}
Soit \(A_5\) le groupe des permutations paires de cinq éléments et soit \(C_5=\langle(1\,2\,3\,4\,5)\rangle\). Ses douze classes à gauche déterminent une représentation par permutations \(\varrho:A_5\to O(\mathbb R^{12})\). La coordinatisation marquée est fixée au moyen des représentants suivants, écrits comme listes d'images :
\[
\begin{array}{c|c@{\qquad}c|c}
j&r_j&j&r_j\\ \hline
1&(4,3,5,2,1)&7&(5,4,3,2,1)\\
2&(4,2,1,3,5)&8&(1,3,4,2,5)\\
3&(1,2,4,5,3)&9&(4,2,3,5,1)\\
4&(2,5,3,4,1)&10&(3,2,5,4,1)\\
5&(4,3,1,5,2)&11&(4,5,2,3,1)\\
6&(5,1,2,3,4)&12&(5,3,2,4,1).
\end{array}
\]
Concrètement, \(\varrho(a)e_j=e_\ell\) si \(ar_jC_5=r_\ell C_5\). Le caractère tridimensionnel \(\chi_3\) prend les valeurs
\[
\begin{array}{c|ccccc}
\text{classe}&1&2^2&3&5_a&5_b\\ \hline
\chi_3&3&-1&0&(1+\sqrt5)/2&(1-\sqrt5)/2.
\end{array}
\]
La classe \(5_a\) contient le générateur de \(C_5\) indiqué et \(5_b\) contient son carré. Ce choix distingue les deux caractères tridimensionnels conjugués. Si \(I_{12}\) est l'identité et \(\mathbf1\) est le vecteur de douze uns, on définit
\begin{equation}
 P_3=\frac3{60}\sum_{a\in A_5}\chi_3(a^{-1})\varrho(a),
 \qquad
 P_{11}=I_{12}-\frac1{12}\mathbf1\mathbf1^{\mathsf T}.
 \label{eq:leech-projectors}
\end{equation}
Le corps \(\mathbb Q(\sqrt5)\) apparaît ici dans la réalisation de la
représentation qui agit sur le registre engendré.

\begin{lemma}[Polarisation de somme nulle]
\label{lem:leech-u}
Les opérateurs de \eqref{eq:leech-projectors} satisfont
\[
 P_3=P_3^{\mathsf T}=P_3^2,\qquad
 \operatorname{rank}P_3=3,\qquad P_3\mathbf1=0,\qquad
 P_3P_{11}=P_{11}P_3=P_3.
\]
Par conséquent,
\begin{equation}
 u=P_3P_{11}K=P_3K,\qquad
 \sum_{j=1}^{12}u_j=0,\qquad
 \|u\|^2=\frac{6638585+2275584\sqrt5}{20}>0.
 \label{eq:leech-u-norm}
\end{equation}
\end{lemma}
\begin{proof}
La somme du caractère est le projecteur isotypique de la représentation.
Le caractère réel et la substitution \(a\mapsto a^{-1}\)
donnent sa symétrie. La multiplicité du caractère tridimensionnel dans la
représentation de \(A_5/C_5\) est la dimension de ses invariants par
\(C_5\), qui vaut
\[
 \frac15\left(
 3+2\frac{1+\sqrt5}{2}+2\frac{1-\sqrt5}{2}
 \right)=1.
\]
Le projecteur a donc rang trois. Son orthogonalité au caractère
trivial implique \(P_3\mathbf1=0\) et les identités avec \(P_{11}\).
L’application de la somme finie au registre
\eqref{eq:leech-K} produit les douze coordonnées
\begin{equation}
\begin{aligned}
u={}&(-495/4-233\sqrt5/10,\quad403/4+169\sqrt5/10,\\
&-403/4-169\sqrt5/10,\quad495/4+233\sqrt5/10,\\
&601/4+1031\sqrt5/10,\quad203/4+211\sqrt5/5,\\
&-203/4-211\sqrt5/5,\quad-601/4-1031\sqrt5/10,\\
&313/4+569\sqrt5/10,\quad162+219\sqrt5/20,\\
&-162-219\sqrt5/20,\quad-313/4-569\sqrt5/10).
\end{aligned}
\label{eq:leech-u-coordinates}
\end{equation}
Les additionner et les élever au carré dans \(\mathbb Q(\sqrt5)\) donne \eqref{eq:leech-u-norm}. De manière équivalente,
\[
 \|u\|^2=K^{\mathsf T}P_3K
 =\frac1{20}\sum_{a\in A_5}
          \chi_3(a^{-1})K^{\mathsf T}\varrho(a)K.
\]
Les coefficients positifs de son évaluation prouvent que \(u\ne0\).
\end{proof}

La somme nulle a une fonction géométrique concrète : les expansions de certains plans pourront compenser les contractions des autres dans le déterminant total. La non-nullité de \(u\), en revanche, garantit que cette compensation contient une déformation effective.

L’espace euclidien ambiant de la réalisation réticulaire est
\begin{equation}
 E=\bigoplus_{j=1}^{12}E_j,\qquad
 B=\operatorname{diag}(B_2,\ldots,B_2),\qquad
 B_2=\begin{pmatrix}2&-1\\-1&2\end{pmatrix}.
 \label{eq:leech-ambient}
\end{equation}
Chaque \(E_j\) est le plan réel engendré par le réseau de racines
\(A_2\), exprimé dans sa base de racines simples. La décomposition
orthogonale appartient à l’espace ambiant. Le réseau de Leech
s’obtient par recollement et passage à un voisin, opérations qui relient
les douze facteurs.

Soit \(C_W\subset\mathbb F_3^{12}\) le code ternaire de dimension six
de la construction de Paley--Witt. En identifiant le discriminant de chaque
facteur \(A_2\) à \(\mathbb F_3\), le recollement est
\[
 N=N(C_W)=
 \bigcup_{c\in C_W}\bigl(A_2^{12}+\phi(c)\bigr),
\]
où \(\phi(c)\) choisit des représentants des classes discriminantes.
L’auto-orthogonalité et l’autodualité du code transmettent parité et
intégralité à \(N\) ; son déterminant est
\[
 \det N=\frac{3^{12}}{(3^6)^2}=1.
\]
La construction exceptionnelle fixe l'origine marquée \(j=1\) et le vecteur
\begin{equation}
 v_\alpha=4\rho_1+\rho_2+\cdots+\rho_{12},
 \qquad \rho_j=(1,1)_j,\qquad
 \|v_\alpha\|_B^2=54.
 \label{eq:leech-marked-vector}
\end{equation}
Avec lui, on forme
\begin{equation}
\begin{split}
 N_{\alpha,0}
  &=\{x\in N:\langle x,v_\alpha\rangle_B\equiv0\pmod3\},\\
 \Lambda&=N_{\alpha,0}+\mathbb Z\,v_\alpha/3.
\end{split}
\label{eq:leech-neighbour}
\end{equation}
La réalisation exceptionnelle établit que \(\Lambda\) est un réseau
positif, pair, unimodulaire, sans racines et de rang \(24\), de sorte qu’il est
isométrique au réseau de Leech \cite{HMT-IV}. En particulier,
\(\Lambda^*=\Lambda\), avec le dual relativement à \(B\).
Ces propriétés et l’appartenance \(v_\alpha/3\in\Lambda\) sont les
hypothèses réticulaires précises utilisées dans les théorèmes suivants.

\subsection{Isométrie ternaire et stabilité du voisin}
\label{subsec:leech-c3}

Nous définissons les matrices bidimensionnelles
\[
 C=\begin{pmatrix}0&-1\\1&-1\end{pmatrix},
 \qquad
 R=\begin{pmatrix}0&1\\1&0\end{pmatrix},
 \qquad
 \varpi=\frac13\begin{pmatrix}2\\1\end{pmatrix}.
\]
La multiplication directe vérifie
\[
 C^3=I_2,\quad C^2+C+I_2=0,\quad
 C^{\mathsf T}B_2C=B_2,\quad
 RCR=C^{-1},\quad R^{\mathsf T}B_2R=B_2.
\]
En outre, \(C\varpi-\varpi=(-1,0)^{\mathsf T}\), de sorte que \(C\)
agit trivialement sur \(A_2^*/A_2\). La réflexion \(R\) change le signe
de la classe discriminante. Le code \(C_W\) contient le mot de
support complet
\[
 c_*=(1,1,1,1,1,1,2,1,1,1,1,1).
\]
Dans la présentation \(C_W=\{(q,qA_W):q\in\mathbb F_3^6\}\), ce mot provient de \(q_*=(1,1,1,1,1,1)\) et de \(q_*A_W=(2,1,1,1,1,1)\).

Pour chaque plan, on fixe le repère
\begin{equation}
 Q_j=\begin{cases}
 R,&(c_*)_j=1,\\
 I_2,&(c_*)_j=2,
 \end{cases}
 \qquad
 g_c=\bigoplus_{j=1}^{12}Q_jCQ_j^{-1}.
 \label{eq:leech-c3}
\end{equation}
Les symboles \(Q_j\) désignent des repères de plans ; le symbole \(P_3\) reste réservé au projecteur sur le secteur tridimensionnel.

\begin{proposition}[Symétrie d’ordre trois du réseau marqué]
\label{prop:leech-c3}
La transformation \(g_c\) est une isométrie d’ordre trois de \(E\),
a pour unique vecteur fixe le vecteur nul et satisfait
\[
 g_cN=N,\qquad g_cN_{\alpha,0}=N_{\alpha,0},
 \qquad g_c\Lambda=\Lambda.
\]
\end{proposition}
\begin{proof}
Chaque bloc est \(C\) ou \(C^{-1}\), dont le polynôme minimal est \(X^2+X+1\). Cela prouve l'ordre et l'affirmation sur les vecteurs fixes, ainsi que la conservation de \(B\). Chaque bloc agit trivialement sur le discriminant et conserve donc toutes les classes de recollement de \(N\).

Écrivons \(v=v_\alpha=\sum_j a_j\rho_j\), avec \(a_1=4\) et
\(a_j=1\) pour \(j>1\). Tous les \(a_j\) sont congrus à un
modulo trois. Les identités
\[
 (C-I_2)\rho=(-2,-1)^{\mathsf T},
 \qquad (C^{-1}-I_2)\rho=(-1,-2)^{\mathsf T}
\]
montrent que
\[
 y=\frac{g_cv-v}{3},\qquad z=\frac{g_c^{-1}v-v}{3}
\]
ont pour classes discriminantes \(c_*\) et \(-c_*\), respectivement.
Les deux appartiennent à \(C_W\), donc \(y,z\in N\). Dans chaque facteur,
\[
 \left\langle
 \frac{(C^{\pm1}-I_2)a_j\rho}{3},a_j\rho
 \right\rangle_{B_2}=-a_j^2.
\]
Par sommation, \(\langle y,v\rangle_B=\langle z,v\rangle_B=-27\) ;
c’est pourquoi \(y,z\in N_{\alpha,0}\).
Si \(x\in N_{\alpha,0}\), alors
\[
 \langle g_cx,v\rangle_B
  =\langle x,g_c^{-1}v\rangle_B
  =\langle x,v\rangle_B+3\langle x,z\rangle_B
  \equiv0\pmod3.
\]
La même opération avec \(g_c^{-1}\) démontre l’égalité
\(g_cN_{\alpha,0}=N_{\alpha,0}\). Enfin,
\(g_c(v/3)=v/3+y\) conserve la classe qui engendre
\(\Lambda=N_{\alpha,0}+\mathbb Z(v/3)\).
\end{proof}

La symétrie ternaire agit donc sur le réseau
recollé complet. Sa conservation utilise le mot du code et
la congruence du vecteur marqué, en plus de la rotation de chaque
plan. Cette précision permet de transporter la symétrie au flot sans
perdre l’information arithmétique qui a rendu sa construction possible.

\subsection{Flux anisotrope et conservation du covolume}
\label{subsec:leech-flow}

\begin{definition}[Déformation déterminée par la polarisation]
Pour \(s\in\mathbb R\), paramètre sans dimension, on définit
\begin{equation}
 H_u=\bigoplus_{j=1}^{12}u_jI_2,\qquad
 D_s=\exp(sH_u)=\bigoplus_{j=1}^{12}e^{su_j}I_2,\qquad
 \Lambda_s=D_s\Lambda.
 \label{eq:leech-flow}
\end{equation}
La coordonnée \(u_j\) est attribuée au plan \(E_j\) de la même coordinatisation marquée.
\end{definition}

Le facteur positif \(e^{su_j}\) conserve l'orientation de chaque plan. Les douze étiquettes transmettent la correspondance entre la polarisation et la réalisation réticulaire. Cette correspondance fait partie de la structure de départ ; une permutation des étiquettes est traitée par leur transport simultané dans les deux objets.

\begin{theorem}[Covolume et symétrie du flux]
\label{thm:leech-flow}
Pour tous \(s,t\in\mathbb R\), on a
\[
\begin{gathered}
 D_{s+t}=D_sD_t,\qquad D_0=I,\qquad D_s^{-1}=D_{-s},\\
 D_s^\dagger=D_s,\qquad D_sg_c=g_cD_s,\qquad \det D_s=1,
\end{gathered}
\]
où \(A^\dagger=B^{-1}A^{\mathsf T}B\).
Chaque \(\Lambda_s\) est un réseau discret de rang \(24\) et de covolume
un, préservé par l’isométrie \(g_c\) d’ordre trois.
\end{theorem}
\begin{proof}
Les identités de groupe se vérifient dans chaque bloc exponentiel.
Les blocs scalaires sont autoadjoints pour \(B_2\) et commutent
avec les blocs de \(g_c\). La somme nulle de
\eqref{eq:leech-u-norm} donne
\[
 \det D_s=\prod_{j=1}^{12}e^{2su_j}
 =\exp\left(2s\sum_{j=1}^{12}u_j\right)=1.
\]
L’image d’un réseau par un opérateur inversible est un réseau de
même rang ; son covolume est multiplié par la valeur absolue du
déterminant. Le covolume de \(\Lambda\) est un, de sorte que celui
de \(\Lambda_s\) l’est également. La commutation et
la proposition~\ref{prop:leech-c3} donnent
\[
 g_c\Lambda_s=g_cD_s\Lambda=D_sg_c\Lambda=D_s\Lambda.
\]
L'ordre et l'espace des points fixes de \(g_c\) demeurent inchangés.
\end{proof}

\begin{remark}[Paramètre et symétrie conservée]
La direction \(u\) est déterminée par le registre ; \(s\) parcourt une collection de réalisations. Remplacer \(u\) par \(u/\|u\|\) reparamétrise la même collection. Une interprétation physique sélectionnant une valeur particulière de \(s\) ajoute sa règle de réalisation. Le théorème conserve la symétrie explicite \(g_c\) ; le transport d'autres automorphismes requiert leurs relations correspondantes avec \(H_u\).
\end{remark}

\subsection{Dual euclidien et déformation effective}
\label{subsec:leech-dual}

\begin{definition}[Dual euclidien]
Pour un réseau \(L\subset E\), son dual par rapport à \(B\) est
\[
 L^*=\{y\in E:\langle y,x\rangle_B\in\mathbb Z
                         \text{ pour tout }x\in L\}.
\]
\end{definition}

\begin{theorem}[Réciprocité du dual]
\label{thm:leech-dual}
La collection \eqref{eq:leech-flow} satisfait
\begin{equation}
 \Lambda_s^*=\Lambda_{-s}\qquad(s\in\mathbb R).
 \label{eq:leech-dual}
\end{equation}
Son autodualité en un paramètre \(s\) équivaut à
\(D_{2s}\Lambda=\Lambda\).
\end{theorem}
\begin{proof}
Pour un opérateur inversible \(A\),
\[
 (AL)^*=A^{-\dagger}L^*.
\]
En effet, \(\langle y,Ax\rangle_B=\langle A^\dagger y,x\rangle_B\)
est entier pour tout \(x\in L\) précisément lorsque
\(A^\dagger y\in L^*\). En prenant \(A=D_s\), l’autoadjonction,
l’identité \(D_s^{-1}=D_{-s}\) et \(\Lambda^*=\Lambda\) prouvent
\eqref{eq:leech-dual}. L’égalité
\(D_s\Lambda=D_{-s}\Lambda\) équivaut, en multipliant par \(D_s\),
à \(D_{2s}\Lambda=\Lambda\).
\end{proof}

Le covolume contrôle le volume d'un domaine fondamental ; l'autodualité contrôle tous les produits entiers avec le réseau. La première propriété est continue le long de cette collection. La seconde impose la condition arithmétique exacte du théorème précédent. On emploiera donc l'expression \emph{covolume un} pour le flot euclidien et l'on spécifiera l'intégralité lorsque l'unimodularité arithmétique intervient.

\begin{proposition}[Témoin de perte locale d'intégralité]
\label{prop:leech-nonintegral}
Il existe \(\delta>0\) tel que, si \(0<|s|<\delta\), le réseau \(\Lambda_s\) n'est pas entier et n'est pas isométrique au réseau de Leech.
\end{proposition}
\begin{proof}
Le vecteur \(v_\alpha/3\) appartient à \(\Lambda\). Sa norme
transportée est
\begin{equation}
 f(s)=\left\|D_s\frac{v_\alpha}{3}\right\|_B^2
 =\frac29\left(16e^{2su_1}+
                       \sum_{j=2}^{12}e^{2su_j}\right).
 \label{eq:leech-witness}
\end{equation}
Par \(\sum_j u_j=0\) et par la première coordonnée de
\eqref{eq:leech-u-coordinates},
\begin{equation}
\begin{split}
 f(0)&=6,\\
 f'(0)&=\frac49\left(16u_1+\sum_{j=2}^{12}u_j\right)
       =\frac{20}{3}u_1
       =-825-\frac{466}{3}\sqrt5\ne0.
\end{split}
\label{eq:leech-witness-derivative}
\end{equation}
La continuité de \(f'\) permet de choisir un intervalle \((-\delta,\delta)\) dans lequel \(f\) est strictement monotone. En réduisant \(\delta\), on obtient en outre \(11/2<f(s)<13/2\). Pour \(0<|s|<\delta\), il en résulte \(f(s)\ne6\) ; dans cet intervalle, six est le seul entier. Par conséquent, \(\Lambda_s\) contient un vecteur de norme non entière. Un réseau entier a une norme entière pour tous ses vecteurs. Puisque Leech est entier et que la norme est conservée par les isométries, la seconde affirmation est également démontrée.
\end{proof}

La déformation préserve volume et symétrie ternaire tout en
modifiant effectivement la structure métrique arithmétique. Les couples
opposés de coordonnées de \(u\) définissent une permutation de l’espace
ambiant qui échange dilatations et contractions. Sa promotion
en un automorphisme du réseau recollé exigerait de conserver le code
et le voisin. La réciprocité du dual déjà démontrée utilise
uniquement l’autoadjonction du flot et l’autodualité initiale.

\subsection{Convergence et réciprocité des fonctions thêta}
\label{subsec:leech-theta}

Pour \(s\in\mathbb R\) et \(t>0\), nous définissons
\begin{equation}
 \Theta_s(t)=\sum_{\lambda\in\Lambda}
                \exp\bigl(-\pi t\|D_s\lambda\|_B^2\bigr).
 \label{eq:leech-theta-real}
\end{equation}
C’est la somme gaussienne de toutes les longueurs du réseau déformé,
avec leurs multiplicités. Si \(M=\max_j|u_j|\), on a
\begin{equation}
 e^{-2|s|M}\|x\|_B^2
 \leq\|D_sx\|_B^2
 \leq e^{2|s|M}\|x\|_B^2.
 \label{eq:leech-quadratic-bounds}
\end{equation}
Sur un compact de \(\mathbb R\times(0,\infty)\), la borne
inférieure fournit une gaussienne dominante sur le réseau fixe.
La série converge absolument et uniformément. Les dérivées
d’ordre fini ajoutent des facteurs polynomiaux dans les coordonnées de
\(\lambda\) ; ces facteurs restent dominés par une autre gaussienne.
Les dérivations terme à terme sont, par conséquent,
légitimes sur ces compacts.

\begin{theorem}[Reciprocidad gaussiana]
\label{thm:leech-theta}
Avec le volume euclidien associé à \(B\),
\begin{equation}
 \Theta_s(t)=t^{-12}\Theta_{-s}(t^{-1}),
 \qquad s\in\mathbb R,\quad t>0.
 \label{eq:leech-theta-reciprocity}
\end{equation}
En particulier, \(\Theta_s(1)=\Theta_{-s}(1)\).
\end{theorem}
\begin{proof}
Nous fixons la transformée de Fourier
\[
 \widehat f(y)=\int_E
      f(x)e^{-2\pi i\langle x,y\rangle_B}\,dx.
\]
Pour un réseau \(L\subset E\) et
\(f_t(x)=e^{-\pi t\|x\|_B^2}\), nous périodisons
\[
 F_t(x)=\sum_{\ell\in L}f_t(x+\ell).
\]
La convergence normale de la série et de ses dérivées donne une fonction lisse sur \(E/L\). Si \(\xi\in L^*\), son coefficient de Fourier est
\[
\begin{split}
 c_\xi
 &=\frac1{\operatorname{covol}(L)}
   \int_{E/L}F_t(x)e^{-2\pi i\langle x,\xi\rangle_B}\,dx\\
 &=\frac{\widehat f_t(\xi)}{\operatorname{covol}(L)}.
\end{split}
\]
La seconde égalité décompose \(E\) en translations d’un domaine
fondamental ; le caractère vaut un sur \(L\). Le produit
des intégrales gaussiennes unidimensionnelles, dans une base
orthonormée pour \(B\), donne
\[
 \widehat f_t(\xi)=
 t^{-12}\exp\bigl(-\pi\|\xi\|_B^2/t\bigr).
\]
La série de Fourier est absolument convergente. En l’évaluant en
\(x=0\), nous obtenons la sommation de Poisson
\[
 \sum_{\ell\in L}e^{-\pi t\|\ell\|_B^2}
 =\frac{t^{-12}}{\operatorname{covol}(L)}
      \sum_{\xi\in L^*}e^{-\pi\|\xi\|_B^2/t}.
\]
À présent, \(L=\Lambda_s\) a un covolume égal à un et \(L^*=\Lambda_{-s}\) d'après le théorème~\ref{thm:leech-dual}. La substitution prouve \eqref{eq:leech-theta-reciprocity}. Le cas \(t=1\) est immédiat.
\end{proof}

La périodisation précédente est la démonstration gaussienne de Poisson ;
elle peut être comparée à \cite[teorema 2, pp.~10--11]{Elkies}.
La convention de signe positif dans Fourier conduit à la même
transformée de la gaussienne, du fait de sa parité. L’égalité des thêta
porte sur les séries complètes et conserve le domaine analytique de
convergence.

\begin{proposition}[Extension holomorphe]
\label{prop:leech-theta-complex}
Pour \(\tau\) dans le demi-plan supérieur
\(\mathbb H=\{\tau\in\mathbb C:\operatorname{Im}\tau>0\}\), la série
\begin{equation}
 \theta_s(\tau)=\sum_{\lambda\in\Lambda}
             e^{\pi i\tau\|D_s\lambda\|_B^2}
 \label{eq:leech-theta-complex}
\end{equation}
converge normalement et satisfait
\begin{equation}
 \theta_s(-1/\tau)=(-i\tau)^{12}\theta_{-s}(\tau).
 \label{eq:leech-theta-modular}
\end{equation}
\end{proposition}
\begin{proof}
La valeur absolue de chaque terme est \(e^{-\pi\operatorname{Im}\tau\|D_s\lambda\|_B^2}\). Sur les compacts du demi-plan supérieur, la partie imaginaire possède une borne inférieure positive, si bien que \eqref{eq:leech-quadratic-bounds} prouve la convergence normale et l'holomorphie. Pour \(\tau=it\), avec \(t>0\), \eqref{eq:leech-theta-reciprocity} appliquée en \(t^{-1}\) donne \(\theta_s(i/t)=t^{12}\theta_{-s}(it)\). Les deux membres de \eqref{eq:leech-theta-modular} sont holomorphes dans \(\mathbb H\) ; le théorème d'identité prolonge l'égalité depuis l'axe imaginaire positif.
\end{proof}

L'inversion modulaire relie les paramètres \(s\) et \(-s\). La périodicité supplémentaire sous \(\tau\mapsto\tau+1\) utilise la parité des normes. La proposition~\ref{prop:leech-nonintegral} présente des paramètres arbitrairement proches de zéro pour lesquels l'intégralité euclidienne est perdue. La collection conserve donc la réciprocité analytique démontrée, tandis qu'une réalisation par des formes modulaires scalaires ou des algèbres vertex de réseaux incorpore en outre les hypothèses arithmétiques correspondantes. Dans la réalisation exceptionnelle initiale, ces hypothèses soutiennent le prolongement vers les algèbres vertex et le Monstrous Moonshine \cite{HMT-IV}.

\subsection{Forme déployée et autodualité en dimension quarante-huit}
\label{subsec:leech-split}

La perte d’intégralité d’une projection euclidienne conduit à
considérer simultanément les deux projections. Dans l’espace
\(E\oplus E\), on introduit la forme bilinéaire
\begin{equation}
 \mathcal B((x,p),(x',p'))=
       \langle x,p'\rangle_B+\langle p,x'\rangle_B.
 \label{eq:leech-split-form}
\end{equation}
Leurs produits apparient les deux coordonnées. Nous définissons
\begin{equation}
\begin{split}
 \mathcal D_s&=D_s\oplus D_{-s},\\
 \Gamma_s&=\mathcal D_s(\Lambda\oplus\Lambda)\\
 &=\{(D_s\lambda,D_{-s}\mu):\lambda,\mu\in\Lambda\}.
\end{split}
\label{eq:leech-split-lattice}
\end{equation}

\begin{theorem}[Autodualité entière pour la forme déployée]
\label{thm:leech-split}
La forme \(\mathcal B\) a pour signature \((24,24)\).
Pour tout \(s\in\mathbb R\), \(\mathcal D_s\) est une isométrie de
\(\mathcal B\), et \(\Gamma_s\) est un réseau entier, pair et autodual
relativement à cette forme.
\end{theorem}
\begin{proof}
En introduisant
\[
 p_L=\frac{x+p}{\sqrt2},\qquad
 p_R=\frac{x-p}{\sqrt2},
\]
on obtient
\[
 \mathcal B((x,p),(x,p))=\|p_L\|_B^2-\|p_R\|_B^2.
\]
Comme \(B\) est positive en dimension \(24\), la signature est \((24,24)\). Le caractère autoadjoint et la réciprocité des blocs donnent
\[
 \langle D_sx,D_{-s}p'\rangle_B=\langle x,p'\rangle_B,
\]
et de manière analogue pour le second produit croisé. Par conséquent,
\(\mathcal B(\mathcal D_s\xi,\mathcal D_s\eta)=\mathcal B(\xi,\eta)\).

Dans \(\Gamma_0=\Lambda\oplus\Lambda\), tous les produits croisés sont entiers et
\[
 \mathcal B((\lambda,\mu),(\lambda,\mu))
   =2\langle\lambda,\mu\rangle_B\in2\mathbb Z.
\]
Cela prouve l’intégralité et la parité. Si \((x,p)\) appartient au dual
de \(\Gamma_0\) pour \(\mathcal B\), l’appariement avec
\((\lambda,0)\) et \((0,\mu)\) exige respectivement
\(p,x\in\Lambda^*=\Lambda\). Le dual est donc \(\Gamma_0\).
L’isométrie \(\mathcal D_s\) transporte ces trois propriétés
à \(\Gamma_s\).
\end{proof}

La métrique spécifie le contenu de la conclusion. La collection \(\Gamma_s\) conserve l'intégralité pour les produits croisés de \(\mathcal B\), même lorsque \(\Lambda_s\) perd l'intégralité pour \(B\). En tant que réseaux abstraits à forme déployée, les \(\Gamma_s\) sont isométriques. Leurs deux projections euclidiennes, considérées avec les métriques positives respectives, varient avec le paramètre.

\begin{proposition}[Échange des projections réciproques]
\label{prop:leech-split-involution}
L'involution \(\mathcal J(x,p)=(p,x)\) satisfait
\[
 \mathcal J\Gamma_s=\Gamma_{-s},\qquad
 \mathcal J\mathcal D_s\mathcal J=\mathcal D_{-s}.
\]
Dans les coordonnées \((p_L,p_R)\), elle fixe \(p_L\) et change le signe de \(p_R\).
\end{proposition}
\begin{proof}
L'échange transforme \((D_s\lambda,D_{-s}\mu)\) en \((D_{-s}\mu,D_s\lambda)\), qui parcourt \(\Gamma_{-s}\). La conjugaison des deux blocs donne la seconde identité. Enfin, \(x+p\) est conservé et \(x-p\) change de signe.
\end{proof}

L’involution exprime une dualité de la construction réticulaire
complète. Son éventuelle représentation physique conserve, en plus de
la forme déployée, les opérateurs et les données d’identification
propres à cette réalisation.

\subsection{Énergie réticulaire et équivalence unitaire}
\label{subsec:leech-energy}

L'énergie de rayon réciproque \(m^2/R^2+w^2R^2\) échange ses deux termes sous \((m,w,R)\mapsto(w,m,R^{-1})\), avec une conjugaison unitaire qui préserve les domaines opératoriels \cite{HMT-IV}. La réalisation sur les douze plans marqués emploie
\begin{equation}
 \mathcal E_s(\lambda,\mu)=
       \|D_s\lambda\|_B^2+\|D_{-s}\mu\|_B^2,
 \qquad \lambda,\mu\in\Lambda.
 \label{eq:leech-energy}
\end{equation}
L’énergie est positive et vérifie
\(\mathcal E_s(\lambda,\mu)=\mathcal E_{-s}(\mu,\lambda)\).

\begin{theorem}[Opérateur positif et conjugaison des domaines]
\label{thm:leech-energy}
Dans l’espace de Hilbert \(\ell^2(\Lambda\oplus\Lambda)\), soit
\[
 (H_s\psi)(\lambda,\mu)=
       \mathcal E_s(\lambda,\mu)\psi(\lambda,\mu)
\]
de domaine
\begin{equation}
 \mathcal D(H_s)=
 \left\{\psi\in\ell^2(\Lambda\oplus\Lambda):
 \sum_{\lambda,\mu}\mathcal E_s(\lambda,\mu)^2
             |\psi(\lambda,\mu)|^2<\infty\right\}.
 \label{eq:leech-energy-domain}
\end{equation}
Alors \(H_s\) est positif, auto-adjoint et à résolvante compacte.
L’application \(U\psi(\lambda,\mu)=\psi(\mu,\lambda)\) est unitaire et
\begin{equation}
 U\mathcal D(H_s)=\mathcal D(H_{-s}),\qquad
 UH_sU^{-1}=H_{-s}.
 \label{eq:leech-unitary}
\end{equation}
\end{theorem}
\begin{proof}
La multiplication par une fonction réelle sur un ensemble discret,
avec le domaine maximal \eqref{eq:leech-energy-domain}, est
autoadjointe. On peut le vérifier directement au moyen de la base
orthonormée des vecteurs à support unitaire : l’appartenance
au domaine de l’adjoint exige que la suite obtenue en
multipliant par \(\mathcal E_s\) soit de carré sommable, soit exactement
la même condition. En outre,
\[
 \langle\psi,H_s\psi\rangle
   =\sum_{\lambda,\mu}\mathcal E_s(\lambda,\mu)
                         |\psi(\lambda,\mu)|^2\geq0.
\]
Par \eqref{eq:leech-quadratic-bounds},
\[
 \mathcal E_s(\lambda,\mu)\geq
 e^{-2|s|M}\bigl(\|\lambda\|_B^2+\|\mu\|_B^2\bigr).
\]
Un réseau a un nombre fini de points dans tout ensemble borné ; chaque sous-niveau de \(\mathcal E_s\) est donc fini. Les coefficients diagonaux \((1+\mathcal E_s)^{-1}\) tendent vers zéro en dehors d'ensembles finis, de sorte que \((H_s+I)^{-1}\) est limite en norme d'opérateurs de rang fini et est compact.

L’échange des deux indices conserve la norme de
\(\ell^2(\Lambda\oplus\Lambda)\), de sorte que \(U\) est unitaire
et \(U^{-1}=U\). L’égalité des énergies lors de l’inversion de \(s\) transforme
la condition de domaine de \(H_s\) en celle de \(H_{-s}\), et
l’évaluation sur chaque paire \((\lambda,\mu)\) prouve
\eqref{eq:leech-unitary}.
\end{proof}

Trois conservations de domaines distincts sont ainsi articulées. Dans les plans euclidiens, le flux maintient le covolume et la symétrie ternaire et relie chaque réseau au dual du paramètre opposé. Dans l'espace doublé, les produits croisés conservent un réseau entier, pair et autodual. Dans la réalisation opératorielle, l'échange réciproque conserve l'énergie au moyen d'une équivalence unitaire d'opérateurs avec leurs domaines. Les formules thêta expriment la même réciprocité dans la somme convergente des longueurs de tout le réseau.
